\documentclass[11pt,letterpaper]{article}
\usepackage[T1]{fontenc}
\usepackage{lmodern}
\usepackage[margin=1in]{geometry}
\usepackage{amsmath,amssymb,amsthm,mathtools}
\usepackage{microtype,booktabs,longtable,array,enumitem}
\usepackage{xcolor,listings,fancyhdr}
\usepackage[colorlinks=true,linkcolor=blue!45!black,urlcolor=blue!50!black,citecolor=blue!45!black]{hyperref}
\usepackage{bookmark}
\setlength{\parskip}{4pt}
\setlength{\parindent}{0pt}
\setlength{\emergencystretch}{2em}
\setlist{nosep,topsep=4pt}
\newtheorem{theorem}{Theorem}[section]
\newtheorem{lemma}[theorem]{Lemma}
\newtheorem{proposition}[theorem]{Proposition}
\newtheorem{corollary}[theorem]{Corollary}
\theoremstyle{definition}
\newtheorem{definition}[theorem]{Definition}
\theoremstyle{remark}
\newtheorem{remark}[theorem]{Remark}
\newcommand{\IPRA}{\mathrm{iPRA}}
\newcommand{\BC}{\mathcal B}
\newcommand{\CC}{\mathcal C}
\newcommand{\fv}{\operatorname{FV}}
\newcommand{\cl}[2]{\langle #1;\,#2\rangle}
\pagestyle{fancy}
\fancyhf{}
\fancyhead[L]{\small iPRA with parameter witnesses}
\fancyhead[R]{\small 12 November 2026}
\fancyfoot[C]{\thepage}
\setlength{\headheight}{14pt}
\lstset{basicstyle=\ttfamily\footnotesize,breaklines=true,columns=fullflexible,
  keepspaces=true,showstringspaces=false,frame=single,rulecolor=\color{black!20},
  xleftmargin=4pt,xrightmargin=4pt,aboveskip=8pt,belowskip=8pt}
\hypersetup{pdftitle={iPRA with parameter witnesses},pdfauthor={OpenAI, prepared for Vladimir Reshetnikov}}

\begin{document}
\hypersetup{pageanchor=false}
\begin{titlepage}
\vspace*{0.55in}
{\Large Research report\par}
\vspace{0.45in}
{\huge\bfseries iPRA with parameter witnesses\par}
\vspace{0.2in}
{\Large A finite witness interface for a complete\par
intuitionistic arithmetic calculus\par}
\vspace{0.5in}
{\large Prepared for Vladimir Reshetnikov\par}
{12 November 2026\par}
\vspace{0.5in}
\begin{abstract}
The current ConciseIPRA development asks whether its paired introduction
rule needs arbitrary numerical witness terms. We prove that it does not.
First, zero and numerical variables suffice. A protected proof translation
then restricts the variable to the newest ambient parameter. The final
presentation uses only one bit: it selects zero or that parameter, with
zero as the empty-scope fallback. The language remains raw scalar iterator
terms and equations with the paired binder $D$. A clause presentation has
five primitive rule families and no separately postulated logical axioms.

The critical arithmetic lemma is a predecessor construction using only
computation, equality and literal-equation induction. Unconditional
Boolean identities remove a circular appeal to numerical case distinction
under an induction hypothesis. A paired formula then names every term by
a fresh variable, recovering every omitted introduction. The resulting
calculus proves exactly the same sequents as the pinned baseline and
inherits both whole-domain interpretation composites with full
first-order intuitionistic primitive recursive arithmetic.

A matched, standalone, publicly scoped JSON proof checker decreases from
2,956 to 2,947 normalized bytes, from 837 to 833 significant tokens, and
from 977 to 971 AST nodes. This modest implementation reduction accompanies
a substantive restriction of the primitive witness interface. It is not
a claim to beat the repository's 516-byte compact codec, nor a proof of
absolute minimality. The new mathematical result is a written proof;
the accompanying executable checks are not a Lean formalization.
\end{abstract}
\vfill
\small Baseline pinned to ProveIt commit
\texttt{9f9545f5ffd9e0392d4a36639503aca12383450a}.\par
\normalsize
\end{titlepage}

\hypersetup{pageanchor=true}
\pagenumbering{roman}
\tableofcontents
\clearpage
\pagenumbering{arabic}

\section{Result and comparison target}

The proposed theory is the five-family calculus $\BC$ of
Section~\ref{sec:bit-calculus}. Its introduction instruction contains
\emph{one bit instead of an arbitrary term}. All arithmetic terms,
ordinary numerical domains, equality, and the reading of the logical
binder are unchanged. This is a restriction of the primitive rule
interface, not a change to the set of mathematical consequences.

The comparison is the repository at the full commit printed on the title
page, inspected on 12 November 2026. The main README records seventy-nine
research packages and several previously established reductions. In
particular, raw clauses, five rule families, computational equality and
two-case evaluation are already present. They are reused here and are
not new claims. The README's further question 2 asks whether introduction
at zero or at a variable suffices, and identifies the recovery of
predecessor without illicit substitution as an obstruction
\cite{baseline,readme}. The first theorem below answers that question;
the strongest version keeps only zero and the newest variable.

\begin{theorem}[Main result]
\label{thm:main}
For every finite numerical scope, every well-scoped context of native
clauses $\Gamma$, and every native Eq/$D$ formula $A$,
\[
 \Gamma\vdash_{\BC} A
 \quad\Longleftrightarrow\quad
 \Gamma\vdash_{\CC} A,
\]
where $\CC$ is the baseline raw clause calculus with unrestricted
witnesses. Consequently the ordinary realization of $\BC$ is
bi-interpretable and equiconsistent with full first-order $\IPRA$.
Both numerical domains are whole and both interpretation composites
have the identity graph $x=y$.
\end{theorem}

The proof is constructive at the metalevel: a finite source derivation
can be translated to a finite target derivation. It does not establish
that the new certificates are shorter; the naming construction generally
makes them longer. The resource being reduced is the trusted rule
interface and its matched checker, not certificate length or checking time.

There is no language-independent smallest checker without a fixed
encoding and implementation model. Accordingly this report establishes
an explicit strict matched reduction and a mathematical restriction. It
does not infer an absolute lower bound from a collection of prior drafts.

\section{Complete definition of the proposed calculus}
\label{sec:bit-calculus}

\subsection{Numerical terms, formulas and scopes}

Terms have one numerical sort:
\[
 t ::= x\mid 0\mid S(t)\mid I(n,a;u.h).
\]
Only the last child $h$ binds $u$. The counter $n$ and seed $a$ do not
bind anything. Every occurrence of $I$ is a finite term description;
there is no object-language variable for a function or its code.
Substitution is ordinary capture-avoiding substitution. Terms are raw:
no normalization, conversion oracle or evaluator is hidden in equality
of syntax trees.

The numerical scope is a finite ordered list $x_1,\ldots,x_k$ of
available free parameters. Local eigenparameters are appended to this
list. Define
\[
 \ell_k=\begin{cases}x_k,&k>0,\\0,&k=0.\end{cases}
\]
This convention concerns proof syntax; it introduces no new numerical
constant or interpretation-domain restriction.

Formulas and clauses are
\[
 A ::= (s=t)\mid D_{x,y}A,
 \qquad
 \cl{A}{y.B(y)},\quad y\notin\fv(A).
\]
The paired binder binds both $x$ and $y$ in its matrix. A context is a
finite list of reusable clauses. An ordinary temporary assumption $F$
abbreviates the constant clause $\cl{0=0}{v.F}$, with $v$ unused.
An empty context contains no implicit axiom clauses.

The ordinary meanings, used in proving adequacy rather than by the
checker, are
\begin{align*}
 (s=t)^\circ&=(s=t),\\
 (D_{x,y}B)^\circ
   &=\exists x\bigl(B^\circ(x,0)\to\forall y\,B^\circ(x,y)\bigr),\\
 \cl{A}{y.B(y)}^\circ&=A^\circ\to\forall y\,B^\circ(y).
\end{align*}
The designated native contradiction is $0=S0$. Ordinary falsity is
discussed explicitly in Section~\ref{sec:biinterpretation}.

\subsection{Five primitive rule families}

All premises retain the indicated context. These are the entire
primitive inventory.

\paragraph{1. Reflexivity.}
For every in-scope term $t$,
\[
 \frac{}{\Gamma\vdash t=t}.\tag{Refl}
\]

\paragraph{2. Evaluation.}
From equations $s=t$ and $b=r$, infer $r=\Phi$, in exactly either of
these two syntactic cases:
\[
\begin{array}{c|c|c}
 b&s&\Phi\\\hline
 I(0,a;u.h)&a&t\\
 I(Sc,a;u.h)&I(c,a;u.h)&h[t/u]
\end{array}
\qquad
 \frac{\Gamma\vdash s=t\qquad\Gamma\vdash b=r}
      {\Gamma\vdash r=\Phi}.
\tag{Eval}
\]
All equalities in the case conditions are literal syntax equality. There
is no third case for $b=s$. The familiar iterator computation equations
and equality transport are derived, as explained below.

\paragraph{3. Clause use.}
For an in-scope term $t$,
\[
 \frac{\Gamma\vdash A}{\Gamma\vdash B(t)}
 \quad\bigl(\cl{A}{y.B(y)}\in\Gamma\bigr).
 \tag{Use}
\]
Clause specialization retains arbitrary terms. Restricting specialization
as well as introduction is not part of this result.

\paragraph{4. Paired elimination.}
\[
 \frac{\Gamma\vdash D_{x,y}B(x,y)\qquad
 \Gamma,\cl{B(b,0)}{y.B(b,y)}\vdash A}
 {\Gamma\vdash A}.
 \tag{Elim}
\]
Here $b$ is fresh for $\Gamma,A$ and all free parameters of $B$.
It is appended to the numerical scope in the second premise. The
displayed clause is discharged.

\paragraph{5. Introduction with one bit.}
Choose $\beta\in\{0,1\}$ and set
\[
 t_\beta=\begin{cases}0,&\beta=0,\\\ell_k,&\beta=1,\end{cases}
 \qquad
 \eta=\begin{cases}
 z,&\beta=0\text{ and }B\text{ is a literal equation},\\
 0,&\text{otherwise}.
 \end{cases}
\]
Then
\[
 \frac{\Gamma,B(t_\beta,\eta)\vdash B(t_\beta,Sz)}
      {\Gamma\vdash D_{x,y}B(x,y)}.
 \tag{Bit-DI}
\]
The eigenparameter $z$ is fresh for the previous scope, context, witness
and matrix parameters. The witness is selected \emph{before} $z$ is
introduced. In particular $z$ itself can never serve as this witness.
The temporary formula assumption is the constant clause defined above.
At an empty numerical scope, bit 1 means a case introduction at zero;
bit 0 at an equation means progress at zero. The two bits therefore still
have different roles there.

\subsection{What is derived}

There are no further primitive logical axioms, successor axioms,
computation axioms, equality schemas, or induction schemas. Hypothesis
use is Refl at $0$ followed by Use of the constant assumption clause.

Eval with reflexive premises gives
\[
 I(0,a;u.h)=a,\qquad I(Sc,a;u.h)=h[I(c,a;u.h)/u].
\]
The first evaluation case also derives the root composition rule
$s=t,\ s=r\vdash r=t$. To see this explicitly, put
$J=I(0,s;u.u)$. Evaluate $s=r$ against $J=J$ to get $J=r$;
then evaluate $s=t$ against $J=r$ to get $r=t$.
Consequently symmetry and transitivity are derived.

For a term template $h(v)$, let
$J=I(0,s;u.h(u))$ and $K=I(S0,s;u.h(u))$.
Zero evaluation derives $J=t$ from $s=t$; successor evaluation with
$K=K$ derives $K=h(t)$. The same calculation with $s=s$ gives
$K=h(s)$, so root composition gives $h(s)=h(t)$. Thus term
congruence and crossed transport follow without a separate transport
instruction. These are the established two-case evaluation reductions
\cite{rootfree,evallean}; they are repeated to make the rule inventory
self-contained.

The rest of the proof derives closure introduction and elimination,
full formula transport, equation induction, predecessor and arbitrary
witness introduction. The last derivation is the new restriction result.
Ordinary intuitionistic connectives are definable, and their ordinary
rules are recovered through the already established native interpretation.
They are not built into the checker.

\section{Why arbitrary witness terms can be removed}

\subsection{The essential naming formula}

Put
\[
 C(v,a,b)=I(v,a;u.b),\qquad
 T(v)=D_{x,y}\bigl(C(y,v,C(v,0,Sx))=v\bigr).
 \tag{Packet}
\]
The variable $u$ is unused in the definition of $C$. The two useful
rows of the matrix are
\[
 y=0:\quad v=v,
 \qquad y=Sw:\quad C(v,0,Sx)=v,
\]
up to the displayed computation equations.

The zero instance $T(0)$ is introduced with bit 0. Its matrix is an
equation, so this is progress; its successor row reduces to $0=0$,
independently of its local hypothesis. The successor instance $T(Sa)$
is introduced with bit 1 while $a$ is the newest parameter. Its successor
row reduces to $Sa=Sa$. Thus an outer closure proves $T(v)$ for every
term $v$ using clause specialization, which remains unrestricted.

Open $T(St)$ at a fresh $b$. Its zero row is a theorem. Specializing
the retained clause at $1=S0$ gives
\[
 C(1,St,C(St,0,Sb))=St,
\]
and computation gives $Sb=St$. A predecessor term $P$ satisfying
$P(Sv)=v$ then gives $b=t$. Hence any proof needing an arbitrary term
$t$ can use a fresh numerical name $b$, transport along $b=t$, and
discharge that name.

This addresses the concrete test case in the README: the ordinary
existential assertion $\exists x\,(x=S0)$ is recoverable. Name $S0$
by opening $T(S(S0))$, infer $b=S0$ by predecessor, introduce the
existential encoding at the available variable $b$, and discharge the
packet. No substitution of $S0$ for a variable in an unverified
restricted derivation is assumed.

\subsection{The arithmetic dependency that prevents circularity}

It would be circular to obtain predecessor from full $\IPRA$
bi-interpretability and then use predecessor to prove that the restricted
theory has full $\IPRA$ strength. Appendix~\ref{app:pure-equational-predecessor}
therefore supplies an internal predecessor construction from computation,
equality and literal-equation induction alone.

The scalar iterator has only a state argument in its update, so predecessor
is not an immediate computation. The inherited construction uses parity,
a three-state cycle, a ramp, a factor map, a comparator and a saturating
counter. The new dependency refinement is important: numerical case
splits under an induction hypothesis are replaced by unconditional
equations proved in advance.

For example, let $N$ be the parity iterator, $j$ the three-state update,
$M$ its orbit, and $G$ the ramp, as defined in the appendix. Before
inducting on $M(G(n))=N(n)$, prove for $b=N(v)$ the identities
\begin{align*}
 M(g+b+1)&=C\bigl(b,j(M(g)),j(j(M(g)))\bigr),\\
 C\bigl(b,j(b),j(j(b))\bigr)&=Z(b).
\end{align*}
Each is obtained from its two direct computations at $b=0,1$ by
equation induction. The ramp step then uses the induction hypothesis
only as an equality, never as a side assumption in a Boolean case split.
This is the precise repair that makes predecessor available before
unrestricted witness introduction.

\subsection{Why the newest variable needs a protected proof}

With zero and arbitrary variables allowed, injective renaming,
weakening and cut are proved directly by induction on derivations.
Appendix~\ref{sec:restricted-witness} gives this simpler proof in full.

Restricting the variable to the newest parameter changes that situation:
inserting another parameter may make a previously permitted witness
cease to be the newest. Thus one cannot silently shift a target proof
under a binder. The proof for the final bit calculus instead maintains
\emph{hereditary derivability}: it constructs a derivation under every
injective renaming and every allowed extension of the retained context.
Under a new eigenparameter, a premise is regenerated from this invariant.
It is not obtained by a forbidden renaming of an existing target tree.

At the final step of arbitrary witness introduction, the packet opens
the name $b$. It is now the newest parameter, so bit 1 introduces at
exactly that name. The step eigenparameter is added only after this
selection. The protected translation and its direct clausal version
are supplied in the proof appendices. This ordering is also exercised
by the eigenvariable-capture rejection test.

\section{Full intuitionistic bi-interpretability}
\label{sec:biinterpretation}

\subsection{The target is full first-order iPRA}

Here $\IPRA$ means first-order intuitionistic predicate logic with
equality, zero and successor, a function symbol for every fixed finite
primitive-recursive description with its defining equations, the usual
successor separation and injectivity axioms, and contextual induction
for quantifier-free formulas. The claim is not restricted to closed
equations, a classical fragment, or a collection of numerical sentences.

The baseline proves native-to-$\IPRA$ and $\IPRA$-to-native translations
for all formulas and contexts, together with both formula round trips.
The pinned Lean endpoint explicitly states
\texttt{native\_iff\_pra}, \texttt{pra\_iff\_native},
\texttt{native\_roundTrip}, \texttt{pra\_native\_roundTrip} and
\texttt{native\_inconsistent\_iff} \cite{biinterp}. These source files
were inspected; the baseline Lean build was not rerun in this workspace.
Our new proof is not included in that existing formalization.

Theorem~\ref{thm:main} transfers this exact result because it identifies
derivability on the unchanged native formulas, with unchanged numerical
terms and clauses. It does not merely show agreement in the standard
natural-number model.

\subsection{The ordinary realization and falsity}

For clarity, elaborate each fixed iterator body $h(u,\vec p)$ into an
ordinary function symbol $i_h(n,a,\vec p)$, with its two defining
equations. Bodies are introduced in finite stages. Let $T_I$ be ordinary
intuitionistic first-order logic with actual equality, these symbols and
equations, induction for literal equations, numerical coverage
\[
 \forall x\,(x=0\lor\exists y\,x=Sy),
\]
and the falsity bridge $(0=S0)\to\bot$.
The native clauses and $D$ are read by the formulas in
Section~\ref{sec:bit-calculus}. Coverage is recovered by the native
logical interpretation. The falsity bridge is explicit in the ordinary
realization and expresses the reading of the designated contradiction;
it is not asserted to follow from positive iterator equations alone.
Neither is an additional instruction in the five-family checker.
Stating the bridge in $T_I$
prevents a mistaken appeal to a singleton model of purely positive
equations.

The forward arithmetic interpretation sends each scalar iterator to its
primitive-recursive iteration in $\IPRA$. The reverse interpretation
first derives predecessor, a reflecting equality test, and quantifier-free
induction. It then defines pairing with projections and compiles every
fixed primitive recursion. For
\[
 f(0,\vec p)=g(\vec p),\qquad
 f(Sn,\vec p)=h(n,f(n,\vec p),\vec p),
\]
the compiler stores the counter and current value in a pair:
\begin{align*}
 V(n,\vec p)&=I\bigl(n,\pi(0,g^*(\vec p));
 u.\pi(SL(u),h^*(L(u),H(u),\vec p))\bigr),\\
 f^*(n,\vec p)&=H(V(n,\vec p)).
\end{align*}
Induction proves $L(V(n,\vec p))=n$, and the projection equations give
the defining equations for $f^*$. Pairing is a derived construction; it
does not enlarge the proof-checking kernel.

\subsection{Both composites, rather than equiconsistency alone}

For each fixed description, its double translation satisfies the same
defining equations. Equation induction proves uniqueness of an iterator
with a fixed seed and update; the analogous argument proves uniqueness
for primitive recursion. External induction on a finite function
description therefore identifies its double translation internally.

Both domains are given by $x=x$, equality is actual equality, and
$x=y$ defines a provably total bijection preserving every fixed
operation in both composites. The inherited logical macro theorems give
the formula composites, including implication, disjunction, conjunction,
both quantifiers and falsity. The two translations preserve contradiction,
so effective translations of finite proofs establish equiconsistency.
No self-consistency statement is asserted.

The unchanged written bridge, equality-reflection and primitive-recursion
sources are included in \texttt{reference/}, with their license and
pinned provenance. The present report proves the new restriction rather
than claiming those inherited constructions as new.

\section{Executable comparison and validation}
\label{sec:implementation}

\subsection{A matched complete checker}

The final module reads one JSON certificate from standard input. It
validates the closed public formula, term and formula scopes, rule
arities, premises and every eigenvariable boundary. It imports only
Python's standard JSON and system modules. It contains no arithmetic
axiom table, proof database, normalizer, or external theorem prover.
Successful acceptance is exit status zero together with the line
\texttt{ACCEPT}; malformed or incorrect certificates are rejected.
As with an ordinary implementation of a proof calculus, practical
recursion and memory limits can reject a valid certificate that exceeds
the host's resources. The equivalence theorem concerns finite formal
derivations, with no fixed machine resource bound.

The baseline is the repository's existing two-case raw clause checker.
Both files use the same representation, same public-goal validation,
same evaluator, same clause handling and same normalization for size
measurement. The change replaces the term-valued introduction field by
a Boolean selector and removes its recursive term-validity check.

\begin{center}
\begin{tabular}{lrrr}
\toprule
Complete standalone module & Tokens & AST nodes & Normalized bytes\\
\midrule
Pinned two-case control & 837 & 977 & 2,956\\
Bit-witness proposal & 833 & 971 & 2,947\\
Decrease & 4 & 6 & 9\\
\bottomrule
\end{tabular}
\end{center}

The source is normalized with \texttt{ast.unparse} under CPython 3.12.14;
the exact scripts and results are included. Only the module docstring is
removed before normalization, consistently for both programs. Comments, documentation and
file formatting are not being used to manufacture the difference. The
byte decrease is about $0.30\%$; token and AST decreases are about
$0.48\%$ and $0.61\%$. The mathematical interface change is larger than
these implementation percentages suggest because the same generic term
validator remains necessary for clause specialization and other rule data.

This comparison does not beat every coordinate of every implementation
in the repository. Its compact 516-byte program is a different codec
profile; the README also records smaller separate token and AST minima.
We do not compare a fully scoped public checker against a compact codec
as if the contracts were identical. There is no global-minimum claim.

\subsection{The Boolean field is a decoder, not a term}

Canonical certificates use the integers 0 and 1 in the introduction
field. The provided checker defines a total Boolean projection of a JSON
value by Python truthiness: false, null, numerical zero, and empty
strings, arrays or objects select bit 0; other parsed values select bit 1.
This is a serialization choice stated as part of the checker contract.
The field is never inserted as a term. The selected numerical term is
constructed internally as zero or the newest scoped free-variable level.

Thus acceptance of a noncanonical truthy mode value cannot create a
new arithmetic witness. An adapter that insists on literal 0 or 1 can
be placed outside this trusted decoder, but no such adapter is silently
excluded from the reported measurement. The measured module itself is
the full specified checker.

\subsection{What the executable evidence establishes}

The package contains actual certificates for the zero naming packet,
its closure, and its instances at 1 and 2. These are checked by the
final restricted checker. Negative cases exercise ill-scoped terms,
incorrect computations, malformed formulas and illicit use of the step
eigenparameter as the introduction witness. The latter is a material
risk: allowing the witness to depend on the new step parameter can
validate a false paired assertion. The final checker constructs the
witness from the preceding scope, so this dependency is excluded.

An independent scalar evaluator also checks 447 finite equalities from
the arithmetic appendix directly from the iterator definitions,
including $P(0),\ldots,P(7)=0,0,1,2,3,4,5,6$. These computations help
catch indexing mistakes. They do not replace the universal equational
proof, prove soundness on nonstandard models, or mechanically certify
bi-interpretability. The validation logs distinguish executed checks
from written theorems.

\section{Scope of the achievement}

The new result is the finite witness interface and the proof that it
preserves full contextual derivability. The arithmetic term language,
the five-family clausal organization, two-case Eval, and the bridge to
$\IPRA$ are inherited and explicitly attributed. The proof also supplies
the missing equational dependency refinement for predecessor.

The implementation result is a small strict reduction of a matched
complete checker. The artifact does not prove an absolute minimum,
establish a new compact-codec record, or provide a newly compiled Lean
formalization. It provides a concrete simpler rule interface, written
proofs of adequacy and both interpretation composites, executable
certificates, exact source provenance and reproducible measurements.

\appendix
% Input this file after defining theorem, lemma, definition, and proof environments.
% Main document should supply amsmath, amssymb, amsthm, hyperref.
\section{The restricted-witness theorem}
\label{sec:restricted-witness}

This section uses raw iterator terms, with ordinary capture-avoiding
substitution.  No computation is built into the syntax.  Write
\[
 C(v,a,b):=I(v,a;u.b),\qquad 1:=S0,
\]
where the bound variable $u$ does not occur in $b$.  Thus
$C(0,a,b)=a$ and $C(Sv,a,b)=b$ are derivable computation equations.
Formulas are equations, $Q_y A$, and $D_{x,y}A$, with their unchanged
ordinary readings
\[
 Q_yA=A(0)\to\forall y\,A(y),\qquad
 D_{x,y}A=\exists x\bigl(A(x,0)\to\forall y\,A(x,y)\bigr).
\]
The distinction between a variable and the numerical constant $0$ is
syntactic throughout.

\subsection{Seven primitive rule patterns}

Let $\mathcal R$ have Hyp, the two iterator computation equations, the
one-term crossed transport rule
\[
 \frac{\Gamma\vdash s=t\qquad\Gamma\vdash h(s)=r}
      {\Gamma\vdash r=h(t)},
\]
and the usual two eliminations
\[
 \frac{\Gamma\vdash Q_yA(y)\qquad\Gamma\vdash A(0)}
      {\Gamma\vdash A(t)},
 \qquad
 \frac{\Gamma\vdash D_{x,y}B(x,y)\qquad
       \Gamma,Q_yB(b,y)\vdash F}
      {\Gamma\vdash F}.
\]
In the second rule $b$ is fresh in the conclusion and the undischarged
context.  Its one introduction rule is
\[
 \frac{\Gamma,B(t,\eta)\vdash B(t,Sz)}
      {\Gamma\vdash D_{x,y}B(x,y)},
 \qquad
 \eta=\begin{cases}
 z,&B\text{ is a literal equation and }t\text{ is the literal }0,\\
 0,&\text{otherwise},
 \end{cases}
 \tag{RDI}
\]
where $z$ is fresh and the supplied witness $t$ must be either the literal
$0$ or a variable.  In particular, every variable witness selects cases.
There is no separate rule for introducing $Q$, and no general term
substitution rule on proofs.

\begin{lemma}[Structural rules available before completeness]
Context weakening, injective renaming of numerical variables, and cut are
admissible in $\mathcal R$.
\end{lemma}
\begin{proof}
Each assertion follows by induction on the given derivation.  A numerical
renaming sends a variable witness to a variable witness and fixes the
literal $0$; it neither changes the root constructor of a formula nor the
introduction mode.  Renamings are extended under each binder by choosing a
fresh variable.  Context weakening does not change a witness.  In the cut
induction the proof replacing a discharged hypothesis is weakened and
injectively renamed under eigenvariables; all other rule data are
unchanged.  This argument does not assert closure under substitution of an
arbitrary term for a numerical variable.  That property will follow only
after unrestricted introduction has been recovered.
\end{proof}

Reflexivity and the usual equational rules are already derived.  For
example, transport $I(0,a;u.u)=a$ along itself, with the identity template,
to obtain $a=a$.  Applying crossed transport with major premise $s=s$,
minor premise $s=t$, and the identity template gives $t=s$.
Transitivity and congruence follow
from crossed transport and reflexivity.  These constructions need no
logical introduction.

\subsection{Unrestricted case introduction for a closure}

\begin{lemma}[Closure introduction]
Both of the following rules are admissible:
\[
 \frac{\Gamma,A(0)\vdash A(Sz)}{\Gamma\vdash Q_yA(y)},
 \qquad
 \frac{\Gamma,E(z)\vdash E(Sz)}{\Gamma\vdash Q_yE(y)},
 \tag{QC, QP}
\]
where $z$ is fresh and $E$ is a literal equation.  The first rule permits
an arbitrary native formula $A$, including an equation.
\end{lemma}
\begin{proof}
For QP use RDI at witness $0$ on $D_{x,y}E(y)$, with $x$ unused, and
then eliminate this dummy witness.  This leaves exactly $Q_yE(y)$.

For QC first note the fixed theorem
\[
 R:=D_{r,w}(0=0).
\]
RDI at witness $0$ proves $R$: its progressive premise concludes the
reflexive equation $0=0$.  Open $R$ at a fresh variable $b$.  In that
context introduce $D_{x,y}A(y)$ at witness $b$.  This is case mode even
when $A$ is an equation, and the required premise is precisely the given
premise, injectively renamed and weakened.  Eliminate the dummy witness
$x$ to obtain $Q_yA(y)$, and then discharge the opening of $R$; neither
$b$ nor the unused assumption $Q_w(0=0)$ occurs in the result.  Thus QC
does not use a forbidden case introduction of an equation at $0$, even
when the original numerical context is empty.
\end{proof}

Consequently, ordinary induction for a literal equation follows by QP
and closure elimination.  All arguments in the separate predecessor
construction can therefore be performed inside $\mathcal R$ once their
equational derivations have been established.

\subsection{Full formula transport without term substitution on proofs}

\begin{lemma}[Leibniz transport]
For every native formula template $H(v)$, the rule
\[
 \Gamma\vdash s=t,\quad\Gamma\vdash H(s)
 \quad\Longrightarrow\quad\Gamma\vdash H(t)
 \tag{L}
\]
is admissible in $\mathcal R$.
\end{lemma}
\begin{proof}
Induct on the number of formula binders in $H$.  State the induction
hypothesis uniformly for every ambient context and every assignment of
terms to the other free variables.  This uniform statement is essential:
we apply the induction hypothesis to newly formed instances, rather than
substitute terms into previously constructed derivations.

For $H(v)=(a(v)=c(v))$, congruence transports both sides and transitivity
joins the resulting equations.

Suppose $H(v)=Q_yA(v,y)$.  Apply QC to the target closure.  Under its
temporary assumption $A(t,0)$, the reverse equation $t=s$ and the
induction hypothesis give $A(s,0)$.  Eliminate the available closure
$Q_yA(s,y)$ at $Sz$ to obtain $A(s,Sz)$.  The forward equation $s=t$ and
the induction hypothesis give $A(t,Sz)$.  This proves the QC premise.
This explicitly handles the contravariant occurrence of the zero row
in the ordinary meaning of $Q$.

Suppose $H(v)=D_{x,y}B(v,x,y)$.  Open the given $H(s)$ at a fresh
variable $b$, obtaining $Q_yB(s,b,y)$.  Introduce the target $H(t)$ at
the permitted variable witness $b$.  Its case premise assumes
$B(t,b,0)$.  Transport that assumption backwards to $B(s,b,0)$,
eliminate the opened closure at $Sz$, and transport the result forwards
to $B(t,b,Sz)$.  All three formulas have fewer binders than $H$.
RDI at $b$ now proves $H(t)$, and the outer elimination discharges $b$.

The equality premise and other existing derivations are only weakened
or injectively renamed when entering a binder.  No application of an
unproved term-substitution principle occurs.
\end{proof}

\subsection{Naming every term with a fresh variable}

Assume the separately proved equational predecessor lemma: there is a
raw iterator term $p(v)$ such that
\[
 \mathcal R\vdash p(Sv)=v.
 \tag{P}
\]
Then successor injectivity is derived by congruence with $p$ followed by
the two instances of (P).  This step must be justified before the
following argument; importing predecessor from the full iPRA
bi-interpretation would be circular.

Put
\[
 A_v(x,y):=\bigl(C(y,v,C(v,0,Sx))=v\bigr),\qquad
 T(v):=D_{x,y}A_v(x,y).
\]
The defining rows have the derived equations
\[
 A_v(x,0)\;\longleftrightarrow\;v=v,
 \qquad
 A_v(x,Sw)\;\longleftrightarrow\;C(v,0,Sx)=v.
\]
Here and below a displayed equivalence between equations abbreviates
transport along the indicated computation equations.

\begin{lemma}[The naming packet]
For every term $t$, $\mathcal R$ proves $T(t)$.  Moreover, the following
fresh-variable naming rule is admissible:
\[
 \frac{\Gamma,b=t\vdash F}{\Gamma\vdash F},
 \qquad b\notin\mathrm{FV}(\Gamma,t,F).
 \tag{Name}
\]
\end{lemma}
\begin{proof}
First prove $T(0)$ by RDI at witness $0$.  Its matrix is an equation, so
the premise is progressive.  Its successor row reduces by the two
selector computations to $0=0$, independently of the induction
hypothesis.

Next, with $z$ a parameter, prove $T(Sz)$ by RDI at the permitted
variable witness $z$.  Its successor row reduces to
$C(Sz,0,Sz)=Sz$, and then to $Sz=Sz$, independently of its case
assumption.  QC applied to $v\mapsto T(v)$ now gives $Q_vT(v)$.
Closure elimination at the arbitrary term $t$, using $T(0)$, gives
$T(t)$.  This last step is an application of QE, not substitution into
the derivation of $T(v)$.

To prove Name, use the theorem $T(St)$ and open its witness at a fresh
$b$.  The resulting assumption is $Q_yA_{St}(b,y)$.  Its zero row is
a computation theorem, so specialization at $1$ gives
\[
 C(1,St,C(St,0,Sb))=St.
\]
The left side computes first to $C(St,0,Sb)$ and then to $Sb$.
Hence $Sb=St$, and (P) gives $b=t$.  Cut this equation into the given
derivation of $F$, and discharge the opened witness.  The conclusion
does not contain $b$.
\end{proof}

\subsection{Recovery of every omitted introduction}

\begin{theorem}[Restriction of witnesses to zero and variables]
Subject to the equational predecessor lemma (P), $\mathcal R$ and the
native scalar kernel derive exactly the same sequents.  In particular,
the combined witness-mode introduction may be restricted to witnesses
that are either $0$ or variables without changing the theory.
\end{theorem}
\begin{proof}
It suffices to recover the native case introduction at an arbitrary term
$t$; native equational progress was recovered by QP.  Suppose
\[
 \Gamma,B(t,0)\vdash B(t,Sz)
 \tag{*}
\]
with the usual freshness conditions.  Name $t$ by a fresh variable $b$.
Inside the context $\Gamma,b=t$, introduce $D_{x,y}B(x,y)$ at the
allowed witness $b$.  Under its temporary assumption $B(b,0)$, full
formula transport gives $B(t,0)$.  Cut this into a weakened and
injectively renamed copy of (*) to obtain $B(t,Sz)$.  Full formula
transport along $t=b$ gives $B(b,Sz)$.  RDI at $b$ proves the desired
$D$ formula.  Name then removes $b=t$ and the fresh variable $b$.

This argument also recovers the omitted case introduction at literal
$0$ for an equation.  Conversely, every rule of $\mathcal R$ is
derived in the native kernel: case mode is its case introduction,
and progressive mode is obtained from its equation-progress closure
followed by case introduction.  Induction on derivations in either
direction completes the sequent equivalence.
\end{proof}

\begin{corollary}[Bi-interpretability and consistency strength]
The ordinary realization of $\mathcal R$ is bi-interpretable with full
first-order iPRA by the same interpretations as the native scalar
kernel.  Both number domains remain whole, equality remains actual
equality, and both domain composites are the identity graph.  In
particular the theories are equiconsistent.
\end{corollary}
\begin{proof}
The syntax and its ordinary reading have not changed, and the theorem
gives equality of all contextual derivability relations, not merely
equality of closed or arithmetically restricted consequences.  Therefore
the already established native bi-interpretation transfers unchanged.
\end{proof}

\paragraph{Relation to the pinned baseline.}
The README at the pinned commit
\texttt{9f9545f5ffd9},
under ``Further questions'', item 2, conjectures the sufficiency of
case-introduction witnesses $0$ and variables, and correctly identifies
the danger of silently using substitution into derivations of the
restricted calculus.  The proof here resolves that particular logical
gap through QC, the explicit packet $T$, and the independently
equational predecessor lemma.  It does not assert any absolute
minimality of the resulting checker, and its mathematical completeness
claim is separate from any program-size measurement.

\paragraph{Auditable source interfaces.}
The pinned baseline's \texttt{NativeEq.lean} establishes the ordinary
one-term equality interface and full formula transport;
\texttt{NativeGate.lean} and \texttt{NativeLogic.lean} establish the
standard logical encodings; and \texttt{NativeBiInterpretation.lean}
states the native/iPRA comparison.  The argument above reproves the
transport and closure steps in the restricted calculus instead of
importing them with their unrestricted proof-substitution lemmas.
The predecessor construction must likewise be given an independent
equational proof; the statement (P) is displayed to make this
dependency explicit.

\section{Restriction to zero or the newest parameter}
\label{sec:newest-witness}

Let $\mathcal R_{\mathrm{new}}$ be the restricted calculus of
Section~\ref{sec:restricted-witness}, with variable witnesses further
restricted to the newest numerical parameter of the current scope.
The literal $0$ remains available.  At an empty scope only $0$ is
available.  Introduction opens its step variable \emph{after} selecting
the witness, so the newest parameter here means the newest parameter
before that step variable is opened.

This restriction is not invariant under arbitrary weakening of a
derivation by a new numerical parameter.  Accordingly, the direct
structural proof from the preceding section cannot be reused.  The
following argument uses a protected translation instead.  It assumes
that the predecessor identity has been proved in the equational
computation/progress fragment, whose derivations are invariant under
numerical renaming; this is stronger than merely assuming an arbitrary
derivation of that identity in the target calculus.

\subsection{Hereditary derivability and uniform cut}

Write $\Gamma\Vdash A$ if, for every injective renaming $\rho$ of the
numerical parameters into any larger numerical scope and every
formula context $\Delta$ containing $\rho\Gamma$, there is an
$\mathcal R_{\mathrm{new}}$ derivation of $\rho A$ from $\Delta$.
This is a metatheoretic family of actual target derivations, not an
additional rule or a second object language.  It is closed by definition
under injective renaming and formula-context weakening.  Membership
$A\in\Gamma$ implies $\Gamma\Vdash A$.

\begin{lemma}[Uniform cut]
If $\Pi$ is an actual target derivation of $A$ from $\Gamma$, and
$\Delta\Vdash B$ for every $B\in\Gamma$ in the same numerical scope,
then there is an actual target derivation of $A$ from $\Delta$.
\end{lemma}
\begin{proof}
Induct on $\Pi$.  At a hypothesis use its given family at the identity
renaming.  At an eigenvariable rule, each replacement assumption has
a new derivation in the enlarged numerical scope because its proof is
hereditary.  Add the temporary assumption as a hypothesis and apply the
induction hypothesis to the subderivation.  The numerical scope and the
chosen witness of the original target rule have not changed, so a
newest-variable witness is still admissible.  No existing target proof
is renamed under the new eigenvariable.  Contexts are scoped: $\Delta$
has the same numerical scope as $\Gamma$, so it cannot already mention
the new eigenparameter of $\Pi$.  Thus its freshness condition is
preserved as well.
\end{proof}

As a consequence, cut is admissible for $\Vdash$: if
$\Gamma\Vdash A$ and $\Gamma,A\Vdash B$, then $\Gamma\Vdash B$.
At each renamed extension, apply uniform cut to the second derivation
using the hereditary family for the first.  More generally a finite
list of hereditary replacement assumptions can be cut simultaneously.

\subsection{The logical and equational interfaces at the hereditary level}

Computation, crossed transport, and QE preserve $\Vdash$.  They introduce
no new numerical parameter, and their conclusion is obtained by applying
the same primitive rule to the available members of the proof families.
Native DE also preserves $\Vdash$: at each renamed extension open the
witness at the newest fresh parameter, and instantiate the hereditary
family for its second premise at the correspondingly extended renaming.
The original eigenparameter maps to this newly chosen parameter.

The closure rules QC and QP hold at the hereditary level.  For QC,
expand the proof through the neutral theorem $D_{r,w}(0=0)$ as in the
preceding section.  Its fresh witness $b$ is newest, and hence is
admissible for the dummy-witness case introduction.  The given step
family is instantiated in the numerical scope containing $b$ and a
new step parameter $z$, rather than obtained by renaming a single
target derivation.  QP uses the literal witness $0$ and is unchanged.

Full formula transport holds for hereditary proofs as well.  The
induction on the formula is exactly the explicit induction in the
preceding section, with hereditary cut in place of ordinary cut.
In the $D$ case the witness that is reintroduced is the one just
opened, so it is newest.  In the $Q$ case use hereditary QC.
All equality and closure premises required below another binder are
regenerated from their hereditary families.  Thus
\[
 \Gamma\Vdash s=t,\quad\Gamma\Vdash H(s)
 \quad\Longrightarrow\quad\Gamma\Vdash H(t).
 \tag{HL}
\]
This statement is uniform in all actual term assignments to free
slots of $H$, as in the previous transport proof.

\subsection{A hereditary naming packet}

Retain $T(v)=D_{x,y}(C(y,v,C(v,0,Sx))=v)$.
The theorem $T(0)$ has a derivation in every numerical scope by
progressive introduction at $0$.
There is also a derivation of $Q_vT(v)$ in every numerical scope.
To see this, \emph{expand} QC instead of invoking it on a supposedly
hereditary proof of $T(Sz)$: open the neutral theorem at a fresh $b$;
introduce the dummy-witness $D$ at $b$; and in its step premise, with
$z$ newly opened, prove $T(Sz)$ by case introduction at $z$.
Here $z$ really is newest, even though $b$ has already been added.
The successor row is a computation theorem.  Eliminate the two dummy
witnesses exactly as before.

This construction works afresh in every ambient numerical scope and
formula context.  Consequently $\varnothing\Vdash Q_vT(v)$ and
$\varnothing\Vdash T(0)$.  QE at an arbitrary actual term $t$ then
gives
\[
 \varnothing\Vdash T(t).
\]
No claim that $T(Sz)$ has a hereditary direct-introduction proof was
needed; such a claim would incorrectly allow $z$ to cease being newest.

Opening $T(St)$ at a fresh newest $b$ gives $b=t$ by computations and
the predecessor identity.  In the opened context, that equality has
a hereditary proof: its derivation uses only the retained closure,
QE, and equational derivations in the renaming-stable fragment.
Therefore it can replace an assumed $b=t$ in an arbitrary target
subderivation by uniform cut.  The packet's DE then discharges $b$.

\subsection{Translation of the native kernel}

\begin{theorem}
The native scalar kernel and $\mathcal R_{\mathrm{new}}$ derive
exactly the same sequents.  Thus in the combined witness-mode rule it
is sufficient to permit the literal $0$ and the newest numerical
parameter, and no other witness.
\end{theorem}
\begin{proof}
Induct on a native-kernel derivation to prove the stronger assertion
$\Gamma\Vdash A$.  All cases except native case introduction follow
from the hereditary interfaces just established.

Consider a case introduction at an arbitrary term $t$, whose premise
is $\Gamma,B(t,0)\vdash B(t,Sz)$.  Its induction hypothesis is a
hereditary family.  To prove the conclusion at a chosen renaming
$\rho$ and context extension $\Delta$, open the hereditary packet
$T(S\rho t)$ at the newest fresh parameter $b$.  It supplies the
hereditarily derivable equation $b=\rho t$.

Temporarily retain $b=\rho t$ as an assumption in addition to
$\Delta$ weakened to the scope containing $b$.  Construct a target
case introduction of the desired $D$ formula at this newest witness
$b$.  In its premise, under the temporary assumption $B(b,0)$ and a
newly opened step parameter, hereditary
transport gives $B(\rho t,0)$.  Instantiate the source premise's
hereditary family at the injection sending the source step parameter
to this new step parameter.  Hereditary cut supplies its missing
zero-row assumption, yielding $B(\rho t,Sz)$.  Hereditary transport
backwards gives $B(b,Sz)$.  Take the member of this step family at the
identity renaming and apply the primitive target introduction at $b$.
Finally use uniform cut to replace the retained assumption
$b=\rho t$ by its hereditary proof from the opened packet; the other
assumptions are supplied by membership in the opened context.  Then
discharge the packet.  The resulting formula is independent of
$b$.

The construction was performed for every $\rho,\Delta$, so it proves
the required hereditary conclusion.  The reverse implication is
immediate: every target primitive rule is already a derived rule of
the native kernel.  Taking the identity member of the hereditary
family completes the comparison.
\end{proof}

The result removes the need to supply an arbitrary term or a variable
index as an introduction witness.  A checker need only choose between
two fixed witnesses, $0$ and the newest parameter, rejecting the
second choice at an empty scope.  This is a statement about the
available witness data, not by itself a claim that a particular
serialization or complete checker is smaller.

\subsection{A total two-choice interface}

For an implementation it is simpler to make both choices total.  Let
$\beta\in\{0,1\}$ be the only introduction datum, and in a numerical
scope of size $n$ put
\[
 t_{n,\beta}:=
 \begin{cases}
 0,&\beta=0\text{ or }n=0,\\
 b_{\mathrm{new}},&\beta=1\text{ and }n>0,
 \end{cases}
 \qquad
 \eta_{B,\beta}:=
 \begin{cases}
 z,&B\text{ is an equation and }\beta=0,\\
 0,&\text{otherwise}.
 \end{cases}
\]
The introduction is the single schema
\[
 \frac{\Gamma,B(t_{n,\beta},\eta_{B,\beta})
        \vdash B(t_{n,\beta},Sz)}
      {\Gamma\vdash D_{x,y}B(x,y)}.
 \tag{Bit-DI}
\]
The mode is selected by the bit, not by inspecting the chosen term.
Thus the choice $\beta=1$ in an empty scope is case introduction at
$0$, including when the matrix is an equation.

\begin{corollary}[One-bit introduction]
Replacing RDI by Bit-DI gives exactly the native scalar kernel's
derivable sequents.
\end{corollary}
\begin{proof}
Every rule of $\mathcal R_{\mathrm{new}}$ is an instance of Bit-DI.
Every instance of Bit-DI is either ordinary case introduction or
equational progress at witness $0$, hence is derived in the native
kernel.  The comparison theorem sandwiches the new calculus between
two calculi with the same derivability relation.
\end{proof}

At the certificate level one may require a literal Boolean datum, or
define a total projection of an arbitrary JSON value onto a bit.
Either is a serialization choice that must be stated explicitly.
The mathematical certificate field no longer represents a term or a
variable index.  In particular its validation never needs to traverse
a proposed witness term.  Replacing the empty-scope rejection by the
specified zero fallback is sound by the corollary; it is not an
unchecked exception to the syntax.

\section{Direct transfer to the five-rule clause checker}
\label{sec:clausal-transfer}

The delivered checker uses the raw Eq/$D$ clause presentation, rather
than formulas with a primitive $Q$.  The newest-witness restriction must
therefore be proved directly for that presentation.  An unrestricted
native-to-clause translation can insert numerical parameters, so it
cannot simply be assumed to preserve the newest-witness condition.

\subsection{The clause interface}

Formulas are equations and $D_{x,y}B(x,y)$.  A retained clause
\[
 [A;y.B(y)]
\]
has ordinary reading $A\to\forall y B(y)$.  An ordinary assumption
$F$ is represented by $\mathsf a(F):=[0=0;y.F]$, with $y$ unused.
The five primitive families are Refl, the two-case Eval, Use, Elim,
and Bit-Intro.  Use takes a proof of the antecedent of a retained
clause and returns its body at a supplied term.  Elim opens a proved
$D_{x,y}B$ at a fresh parameter $b$ by retaining
\[
 [B(b,0);y.B(b,y)]
\]
in its branch.  Bit-Intro adds the ordinary-assumption clause
$\mathsf a(B(t_{n,\beta},\eta_{B,\beta}))$ in its step branch and
checks the conclusion $B(t_{n,\beta},Sz)$, with the selectors from
Bit-DI.  All terms are raw; no normalization is implicit.

Refl and Eval alone derive both iterator computations and the whole
equational interface.  For example, the zero case of Eval is
\[
 s=t,\quad I(0,s;u.h)=r\quad\Longrightarrow\quad r=t.
\]
Two uses, through the reflexive frame $I(0,s;u.h)$, derive the root
joining rule $s=t,\ s=r\Longrightarrow r=t$.  The successor case
applied to a one-step iterator then gives congruence in an arbitrary
term template.  All these derivations are uniform in the numerical
scope and use no introduction or elimination of $D$.  They therefore
remain valid when regenerated under any numerical renaming.

Define the closure abbreviation
\[
 q_yA(y):=D_{x,y}A(y),\qquad x\text{ unused}.
\]
For the total bit interface, case introduction of $qA$ is a direct
Bit-Intro with $\beta=1$: the witness is unused, and the empty-scope
fallback still selects cases.  Equational progress introduction of
$qE$ is direct Bit-Intro with $\beta=0$.  Closure elimination is an
Elim of $qA$, followed by Use of the resulting clause
$[A(0);y.A(y)]$.  These observations implement the $Q$ interface
without a primitive $Q$ constructor.  Hereditary premises, defined
next, supply any proofs needed below the fresh dummy witness.

\subsection{Hereditary derivations with clause contexts}

Write $K\Vdash_c F$ if every injective numerical renaming of $K,F$
into any larger finite numerical scope is derivable in the bit clause
calculus from every clause-context extension of the renamed $K$.
Renaming acts capture-avoidably on
both the antecedent and the bound body of each clause.  The extension
retains clauses, including their entire antecedent/body data; a
retained clause is not replaced by the assertion that its antecedent
or one instance of its body is true.

The cut principle needed here concerns \emph{ordinary-assumption}
clauses only:
\[
 K,\mathsf a(A)\vdash_c F,\qquad K\Vdash_c A
 \quad\Longrightarrow\quad K\vdash_c F.
 \tag{UC$_c$}
\]
It also holds when the other retained clauses are embedded in a
larger target context.  Prove it by induction on the actual first
derivation.  When Use selects the distinguished clause
$\mathsf a(A)$, replace that use by the available member of the
hereditary family for $A$; its antecedent proof of $0=0$ is no longer
needed.  Every other clause use keeps the same clause and applies
the induction hypothesis to its antecedent proof.  Elim and Intro
add their own clauses in both contexts; the replacement family is
regenerated in the enlarged numerical scope.  Thus no arbitrary
clause is cut or semantically reconstructed, and no newest-sensitive
proof is renamed.  The same argument gives cut for
$\Vdash_c$ at an ordinary-assumption clause.

Refl, Eval, and Use preserve $\Vdash_c$.  For Use the selected clause
is a member of the renamed context, and the hereditary premise
supplies its antecedent.  Elim preserves $\Vdash_c$ by extending
the source injection to send its eigenparameter to the newest
fresh target parameter; its new clause is retained literally in
the target branch.

Full formula transport is proved by induction on Eq/$D$ formulas
at the hereditary level.  The equation case is the equational
interface above.  For $D_{x,y}B(s,x,y)$, open its witness at the
newest $b$ and retain the clause
$[B(s,b,0);y.B(s,b,y)]$.  Introduce the target $D$ at $b$ with bit 1.
From its temporary ordinary assumption $B(t,b,0)$, the induction
hypothesis transports backwards to $B(s,b,0)$.  Use of the retained
clause at $Sz$ gives $B(s,b,Sz)$; forward transport gives
$B(t,b,Sz)$.  Hereditary ordinary cut composes these operations.
The result is independent of $b$, so Elim discharges it.  This proves
\[
 K\Vdash_c s=t,\quad K\Vdash_c F(s)
 \quad\Longrightarrow\quad K\Vdash_c F(t).
 \tag{HL$_c$}
\]
The induction is uniform in actual term assignments to the other
free slots.  No arbitrary term is substituted into a target proof.

\subsection{The packet and the completeness induction}

The same raw equation packet works:
\[
 T(v):=D_{x,y}\bigl(C(y,v,C(v,0,Sx))=v\bigr).
\]
At any ambient scope, $T(0)$ is obtained with bit 0 by a
progressive introduction whose successor row is a computation
theorem.  To construct $q_vT(v)$, use bit 1 on the outer dummy
witness.  In its step branch the parameter $z$ is newest; construct
$T(Sz)$ with bit 1, choosing exactly that $z$.  Its successor row
again computes to reflexivity.  This construction is repeated in
each ambient scope.  It does not assume that a direct proof of
$T(Sz)$ remains valid after $z$ ceases to be newest.

Elim followed by Use of the resulting closure clause gives
$K\Vdash_c T(t)$ for every term $t$.  Opening $T(St)$ at the
newest $b$ and using its clause at $1$ gives $Sb=St$.  The
independently equational predecessor identity gives $b=t$.
That equality has a hereditary proof in the opened clause context,
because the calculation uses only retained-clause Use and the
renaming-stable equational fragment.

Now induct on an unrestricted clause-calculus derivation, proving
hereditary derivability of its conclusion.  Refl, Eval, Use and
Elim were handled above.  A progressive introduction at witness
$0$ is directly bit 0.  For a case introduction at any term $t$,
open $T(St)$ at a newest fresh $b$.  Temporarily retain
$\mathsf a(b=t)$ and introduce the desired $D$ with bit 1.
In its step branch, hereditary transport changes the temporary
assumption $B(b,0)$ into $B(t,0)$.  Instantiate the induction
hypothesis for the original source premise in the enlarged scope:
the original source variables keep their chosen injection, the
source step parameter maps to the newest target step parameter,
and the alias parameter $b$ is skipped by this injection.
Ordinary hereditary cut supplies the zero-row assumption, and
transport backwards gives $B(b,Sz)$.

After this introduction, UC$_c$ replaces the retained ordinary
assumption $b=t$ by its hereditary proof from the opened packet.
Elim discharges $b$.  The argument applies at every numerical
renaming and clause-context extension, so it proves the hereditary
induction claim.

Conversely, every bit introduction is either ordinary cases or
equational progress.  The only extra case relative to the
baseline's witness-tested combined rule is bit 1 at an empty
scope with an equation matrix: this is case introduction at $0$.
It is derived in the unrestricted raw clause calculus by using the
nonliteral alias $w=I(0,0;u.0)$ as its witness and transporting
along $w=0$.  Since $w$ is not the literal $0$, the baseline rule
selects cases.  Thus no semantic strengthening occurs at the
empty-scope fallback.

\begin{theorem}[Five-rule bit clause calculus]
The five-rule raw Eq/$D$ clause calculus with Bit-Intro derives
exactly the same sequents from the same clause contexts as the
unrestricted clause calculus.  In particular it has the same
ordinary realization, and hence the same full first-order
bi-interpretation with iPRA.
\end{theorem}

The proof is a direct transformation into the restricted clause
calculus.  It does not rely on an intermediate proof translation
preserving which numerical parameter happens to be newest.

\section{A predecessor from equations and equation induction alone}
\label{app:pure-equational-predecessor}

This appendix supplies the arithmetic fact needed by the restricted-witness
theorem: successor injectivity follows from a purely equational fragment,
without case introduction for $D$.  The construction of the predecessor is inherited
from the repository's scalar bootstrap.  The contribution of this appendix
is its purely equational organization, in particular the replacement of
case splits made under an induction hypothesis by unconditional finite
identities proved in advance.  The argument below uses no numerical
coverage axiom, no logical disjunction, and no successor injectivity.
After equation induction is made available through the interface below,
the arithmetic argument itself has no $D$ formulas.

\subsection{Resources and a finite-value lemma}

The available terms are variables, $0$, $S$, and the iterator
$I(n,a;u.h)$, whose only bound variable is $u$ in $h$.  We use
\begin{align}
 I(0,a;u.h)&=a, \label{eq:pe-zero}\\
 I(Sn,a;u.h)&=h[I(n,a;u.h)/u], \label{eq:pe-succ}
\end{align}
ordinary equality reasoning, and induction on a single equation.  In the
native calculus the latter is obtained from QI by progress and QE:
\[
 \frac{\Gamma\vdash E(0)\qquad
       \Gamma,E(z)\vdash E(Sz)}{\Gamma\vdash E(t)},
 \qquad z\notin\operatorname{FV}(\Gamma),
\]
where $E$ is a literal equation and all its other parameters stay fixed.
In a combined-introduction presentation, QI by progress is obtained by
introducing an equation matrix with unused witness at $0$ and then opening
that binder.  The equational argument itself needs only the displayed
equation-induction interface.

\paragraph{Instantiation in the equational fragment.}
Capture-avoiding substitution of arbitrary terms for numerical parameters
is admissible in this equational fragment.  This follows by induction on
an equational derivation: computation and transport commute with
substitution, and an induction eigenvariable is first renamed fresh for
the terms being substituted.  In the combined-rule implementation of
equation induction, the witness remains the literal $0$ and the matrix
remains an equation.  Thus this argument never substitutes into a case
introduction at a variable witness.  Completed equational identities may
therefore be instantiated at arbitrary terms without presupposing a
substitution theorem for the full restricted calculus.

Ignoring the induction hypothesis gives \emph{equational cases}: proofs of
$E(0)$ and $E(Sz)$ give $E(t)$ for every term $t$.  This is not a case rule
for arbitrary formulas.  It is a derived rule only for equations.

Write $1=S0$, with larger numerals defined similarly, and put
\[
 C(n,a,b)=I(n,a;u.b),\qquad
 Z(n)=C(n,1,0),\qquad
 N(n)=I(n,0;u.Z(u)).
\]
The update variable is absent from the body of $C$, so computation gives
\begin{gather}
 C(0,a,b)=a,\qquad C(Sn,a,b)=b,\label{eq:pe-selector}\\
 Z(0)=1,\qquad Z(Sn)=0,\label{eq:pe-zero-test}\\
 N(0)=0,\qquad N(Sn)=Z(N(n)).\label{eq:pe-parity-rec}
\end{gather}

\paragraph{Finite-value transfer.}
Let $F(v)$ be an equation, and suppose that $F(0)$ and $F(1)$ have been
proved over a fixed side context $\Gamma$.  Then $F(Z(t))$ and $F(N(t))$
are derivable over $\Gamma$ for every term $t$.  The placeholder $v$ is
fresh for $\Gamma$; all other parameters of $F$ remain fixed.

For the first conclusion, use equational cases on a fresh variable $x$
with equation $F(Z(x))$.  Its zero case is $F(1)$ and its successor case is
$F(0)$, by \eqref{eq:pe-zero-test}.  For the second, use equational cases
on $x$ with equation $F(N(x))$.  Its zero case is $F(0)$; its successor
case is $F(Z(N(x)))$, already provided by the first conclusion.
Finally instantiate the completed equations at $x=t$.

In subsequent uses of this lemma, the two equations $F(0)$ and $F(1)$
are unconditional computations.  We never infer an equation by splitting
on $N(t)=0$ or $N(t)=1$ while a side assumption depends on that value.

\subsection{Addition and parity identities}

Define
\[
 a+b=I(b,a;u.Su),\qquad 2a=a+a,\qquad 3a=a+a+a.
\]
Thus $a+0=a$ and $a+Sb=S(a+b)$.  Equation induction on $b$ gives
$0+b=b$ and $Sa+b=S(a+b)$.  Another induction on $b$ proves
commutativity: the successor step is
\[
 a+Sb=S(a+b)=S(b+a)=Sb+a.
\]
Induction on the last summand proves associativity.  We shall reorder
finite sums using these completed identities.  In particular,
$a+1=Sa$, $2(Sa)=2a+2$, and $3(Sa)=3a+3$.

The instances at $v=0,1$ of
\[
 v+Z(v)=1,\qquad Z(Z(v))=v
\]
compute directly.  Finite-value transfer therefore gives
\begin{equation}
 N(n)+Z(N(n))=1,\qquad Z(Z(N(n)))=N(n).
 \label{eq:pe-parity-bool}
\end{equation}
Two successor computations in \eqref{eq:pe-parity-rec} now give
\begin{equation}
 N(n+2)=N(n).\label{eq:pe-period-two}
\end{equation}
An equation induction on $n$ proves $N(2n)=0$: the successor step reduces
$N(2n+2)$ to $N(2n)$.  Consequently $N(2n+1)=Z(0)=1$.
An independent equation induction proves $N(3n)=N(n)$: after increasing
$n$, the left argument increases by three, so period two reduces the
new left value to $Z(N(3n))$, and the induction equation changes this to
$Z(N(n))=N(Sn)$.  We have proved
\begin{equation}
 N(2n)=0,\qquad N(2n+1)=1,\qquad N(3n)=N(n).
 \label{eq:pe-parity-scaling}
\end{equation}
No assertion that every number is even or odd has been used.

\subsection{A three-state update and a ramp}

Put
\[
 j(u)=Z(u)+2N(u),\qquad
 M(n)=I(n,0;u.j(u)),\qquad
 G(n)=I(n,0;u.u+M(u)+1).
\]
Finite computations give
\begin{equation}
 j(0)=1,\qquad j(1)=2,\qquad j(2)=0.
 \label{eq:pe-three-cycle}
\end{equation}
We need these values and $M(Sn)=j(M(n))$; we do not assume a remainder
theorem for $M$.

For arbitrary parameter $g$, the following equations hold whenever
$b=N(v)$:
\begin{align}
 M(g+b+1)&=C\bigl(b,j(M(g)),j(j(M(g)))\bigr),
    \label{eq:pe-ramp-finite-a}\\
 C\bigl(b,j(b),j(j(b))\bigr)&=Z(b).
    \label{eq:pe-ramp-finite-b}
\end{align}
To prove \eqref{eq:pe-ramp-finite-a}, apply finite-value transfer to the
equation with placeholder $b$.  At $b=0$ it is
$M(g+1)=j(M(g))$, and at $b=1$ it is
$M(g+2)=j(j(M(g)))$, both computation identities.  For
\eqref{eq:pe-ramp-finite-b}, the zero instance is
$j(0)=Z(0)=1$, and the one instance is
$j(j(1))=j(2)=Z(1)=0$.  These are again computations.  These two
identities are proved \emph{before} the next induction.

We claim
\begin{equation}
 M(G(n))=N(n).\label{eq:pe-ramp-invariant}
\end{equation}
The zero instance is immediate.  In the successor step retain the single
induction equation $M(G(n))=N(n)$ and calculate
\begin{align*}
 M(G(Sn))
 &=M\bigl(G(n)+M(G(n))+1\bigr)\\
 &=M\bigl(G(n)+N(n)+1\bigr)\\
 &=C\bigl(N(n),j(M(G(n))),j(j(M(G(n))))\bigr)\\
 &=C\bigl(N(n),j(N(n)),j(j(N(n)))\bigr)\\
 &=Z(N(n))=N(Sn).
\end{align*}
The second and fourth lines use the induction equation; the third and
fifth use the two unconditional identities just proved.  In particular,
there is no case split under the induction hypothesis.

It follows that
\begin{equation}
 G(Sn)=G(n)+N(n)+1,\qquad G(n+2)=G(n)+3.
 \label{eq:pe-ramp-step}
\end{equation}
The second equation follows from the first applied twice and
$N(n)+Z(N(n))=1$.  Equation induction gives
\begin{equation}
 G(2n)=3n.\label{eq:pe-ramp-double}
\end{equation}
For the successor step, use $G(2n+2)=G(2n)+3$.

A separate equation induction proves
\begin{equation}
 G(3n)=3G(n)+N(n).\label{eq:pe-ramp-triple}
\end{equation}
Indeed, \eqref{eq:pe-ramp-step} gives
$G(x+3)=G(x)+N(x)+4$.  Assuming the induction equation and using
$N(3n)=N(n)$, the new left side is
\[
 G(3Sn)=G(3n)+N(3n)+4=3G(n)+2N(n)+4.
\]
The new right side is
\[
 3G(Sn)+N(Sn)=3G(n)+3N(n)+3+Z(N(n)).
\]
The two expressions are equal by $N(n)+Z(N(n))=1$ and the additive
identities.  The base is $0=0$.  Again, no conditional case distinction
occurs inside the induction.

\subsection{The factor map}

Define
\[
 H(u)=C(N(u),G(u),u).
\]
Then
\begin{equation}
 H(2u)=3u,\qquad H(1)=1,\qquad H(3u)=3H(u).
 \label{eq:pe-factor}
\end{equation}
The first two equations follow immediately from
\eqref{eq:pe-parity-scaling}, \eqref{eq:pe-ramp-double}, and the two
selector computations.  For the third, first prove by finite-value
transfer the unconditional equation
\begin{equation}
 C(b,3a+b,3c)=3C(b,a,c),\qquad b=N(v).
 \label{eq:pe-factor-finite}
\end{equation}
At $b=0$ both sides are $3a$, and at $b=1$ both sides are $3c$.
Now substitute $v=u$, $a=G(u)$, and $c=u$ into this completed identity:
\begin{align*}
 H(3u)
 &=C\bigl(N(u),3G(u)+N(u),3u\bigr)\\
 &=3C(N(u),G(u),u)=3H(u).
\end{align*}
This argument does not represent an arbitrary number as $2v$ or $2v+1$.

\subsection{A comparator}

For a fixed update term $h$, whose parameters do not include the variable
being inducted on, equation induction proves the rotation identity
\begin{equation}
 I(Sy,a;h)=I(y,h(a);h).
 \label{eq:pe-rotation}
\end{equation}
At $y=0$ both sides are $h(a)$.  In the step, the two successor
computations apply the same fixed $h$ to the induction equation.
Here $I(y,a;h)$ abbreviates $I(y,a;u.h(u))$; it is a term abbreviation,
not quantification over functions.

Using $H(3u)=3H(u)$, another equation induction on $y$, with $a$ fixed,
gives
\begin{equation}
 I(y,3a;H)=3I(y,a;H).
 \label{eq:pe-orbit-scaling}
\end{equation}
The base is $3a=3a$; the step rewrites
$H(I(y,3a;H))$ to $H(3I(y,a;H))=3H(I(y,a;H))$.
Similarly, $H(1)=1$ gives
\begin{equation}
 I(y,1;H)=1.\label{eq:pe-orbit-one}
\end{equation}

Set
\[
 A(x)=I(x,1;u.2u),\qquad
 L(x,y)=N\bigl(I(y,A(x);H)\bigr).
\]
We obtain the three comparison equations
\begin{equation}
 L(0,y)=1,\qquad L(Sx,0)=0,\qquad L(Sx,Sy)=L(x,y).
 \label{eq:pe-comparator}
\end{equation}
The first uses $A(0)=1$, \eqref{eq:pe-orbit-one}, and $N(1)=1$.
The second uses $A(Sx)=2A(x)$ and $N(2u)=0$.  For the third,
\begin{align*}
 L(Sx,Sy)
 &=N\bigl(I(Sy,2A(x);H)\bigr)\\
 &=N\bigl(I(y,H(2A(x));H)\bigr)\\
 &=N\bigl(I(y,3A(x);H)\bigr)\\
 &=N\bigl(3I(y,A(x);H)\bigr)\\
 &=N\bigl(I(y,A(x);H)\bigr)=L(x,y).
\end{align*}
The successive steps use rotation, $H(2u)=3u$, orbit scaling, and
$N(3u)=N(u)$.  Thus \eqref{eq:pe-comparator} holds internally at all
arguments.  The intended interpretation of $L(x,y)$ as the indicator of
$x\leq y$ is not a premise.

\subsection{Saturation and predecessor}

Define the following open terms, with $m$ a fixed parameter:
\[
 B_m(u)=C(L(Su,m),u,Su),\qquad
 W(t,m)=I(t,0;u.B_m(u)).
\]
Equational cases on a fresh variable $e$ prove
\begin{equation}
 C(e,Sa,Sb)=S C(e,a,b).
 \label{eq:pe-selector-succ}
\end{equation}
At $e=0$ both sides are $Sa$, and at $e=Sz$ both are $Sb$.
Using \eqref{eq:pe-comparator} and this selector equation gives
\begin{align}
 B_{Sm}(Su)
 &=C(L(SSu,Sm),Su,SSu)\notag\\
 &=C(L(Su,m),Su,SSu)=S B_m(u),
 \label{eq:pe-saturation-conjugacy}\\
 B_{Sm}(0)
 &=C(L(S0,Sm),0,S0)\notag\\
 &=C(L(0,m),0,S0)=C(1,0,S0)=S0.
 \label{eq:pe-saturation-base}
\end{align}

Equation induction on $t$, with $m$ fixed, now proves
\begin{equation}
 W(St,Sm)=S W(t,m).
 \label{eq:pe-saturation-shift}
\end{equation}
For $t=0$, the left side is $B_{Sm}(0)=S0$, as required.  The step is
\begin{align*}
 W(SSt,Sm)
 &=B_{Sm}(W(St,Sm))\\
 &=B_{Sm}(S W(t,m))\\
 &=S B_m(W(t,m))=S W(St,m).
\end{align*}
Only the second line uses the active induction equation.

Two more equation inductions give
\begin{equation}
 L(Sn,n)=0,\qquad W(n,n)=n.
 \label{eq:pe-saturation-diagonals}
\end{equation}
For the first, the base is $L(S0,0)=0$ and the step reduces
$L(SSn,Sn)$ to $L(Sn,n)$.  For the second, the base is $W(0,0)=0$,
and the step uses the completed identity
\eqref{eq:pe-saturation-shift} at $t=m=n$:
$W(Sn,Sn)=S W(n,n)=Sn$.  This substitution is into an already proved
identity, and does not change a parameter of an active induction
hypothesis.  Consequently
\begin{equation}
 W(Sn,n)=B_n(n)=C(L(Sn,n),n,Sn)=n.
 \label{eq:pe-saturation-extra}
\end{equation}

Finally define the finite scalar term
\begin{equation}
 P(m)=I\bigl(m,0;u.C(L(SSu,m),u,Su)\bigr).
 \label{eq:pe-predecessor-term}
\end{equation}
Its zero equation is a computation.  For its successor equation put
\[
 V(t,m)=I\bigl(t,0;u.C(L(SSu,Sm),u,Su)\bigr).
\]
The third equation in \eqref{eq:pe-comparator} gives the pointwise
identity
\[
 C(L(SSu,Sm),u,Su)=C(L(Su,m),u,Su)=B_m(u).
\]
Equation induction on $t$, with $m$ fixed, proves $V(t,m)=W(t,m)$.
The base is $0=0$.  In the step, first apply the pointwise update
identity at $u=V(t,m)$ and then the induction equation, obtaining
$B_m(V(t,m))=B_m(W(t,m))$.  Thus no primitive extensionality rule for
iterator bodies is being assumed.  From
\eqref{eq:pe-saturation-extra} we conclude
\begin{equation}
 \boxed{\quad P(0)=0,\qquad
 P(Sm)=V(Sm,m)=W(Sm,m)=m.\quad}
 \label{eq:pe-predecessor-equations}
\end{equation}
Applying equality congruence for $P$ to $Sx=Sy$ and using
\eqref{eq:pe-predecessor-equations} gives $x=y$.

Every abbreviation in this appendix expands to a finite term over
$0,S,I$.  Every induction invariant is one literal equation.  Equality
reflection, pairing, general primitive-recursion compilation, numerical
coverage, and logical case analysis were not used to obtain predecessor
or successor injectivity.

\subsection{Provenance and verification status}

The arithmetic construction, including the factor map and the saturated
iteration, comes from
\texttt{Research/certified\_zero/reference/scalar\_predecessor.tex} in the
user's ConciseIPRA project, read at commit
\texttt{9f9545f5ffd9e0392d4a36639503aca12383450a}.
The repository URL is
\begin{center}\small
\href{https://github.com/VladimirReshetnikov/ProveIt/tree/9f9545f5ffd9e0392d4a36639503aca12383450a/Logic/PrimitiveRecursiveArithmetic/ConciseIPRA}{ConciseIPRA at the pinned commit}
\end{center}
The existing source explicitly allowed reconstructed numerical cases.
The present appendix strengthens its dependency statement by giving
equation-only proofs of the needed finite identities, particularly
\eqref{eq:pe-ramp-finite-a}, \eqref{eq:pe-ramp-finite-b}, and
\eqref{eq:pe-factor-finite}.  These identities remove the circularity that
would arise if a proof of restricted logical introduction already assumed
unrestricted logical case distinction.

This appendix is a written internal derivation.  It does not claim a newly
compiled Lean formalization of the equational dependency refinement or of
the restricted-witness theorem.  The repository's pre-existing Lean
formalization should not be presented as machine verification of this
new dependency claim without the corresponding new formal proof.


\section{Archive map and reproduction}

\begin{longtable}{p{0.40\linewidth}p{0.54\linewidth}}
\toprule
Path & Purpose\\\midrule
\endhead
\texttt{ipra\_parameter\_witness.tex}, PDF & Main report and complete new proof appendices.\\
\texttt{paper/} & LaTeX sources for the restriction, protected translation and arithmetic.\\
\texttt{checker/checker.py} & Standalone five-family bit-witness checker.\\
\texttt{checker/reference/} & Exact pinned unrestricted checker controls.\\
\texttt{checker/reproduce.py} & Exact source metrics and certificate checks.\\
\texttt{code/check\_scalar\_arithmetic.py} & Finite evaluation from iterator definitions.\\
\texttt{checker/fixtures/}, \texttt{evidence/} & Checked certificates, tests, metrics and validation logs.\\
\texttt{reference/} & Attributed baseline bridge sources and inspected Lean endpoints.\\
\texttt{PROVENANCE.json} & Exact repository paths, commit and source hashes.\\
\texttt{CLAIMS.json} & Claims, dependencies and explicit limitations.\\
\texttt{reproduce.py} & Re-run executable validation and measurements.\\
\texttt{build.py} & Rebuild the PDF with installed pdfLaTeX.\\
\bottomrule
\end{longtable}

Run \texttt{python3 reproduce.py} from the extracted directory.
Run \texttt{python3 build.py} to rebuild the PDF; a normal LaTeX
installation with the packages named in the source is required.
The executable checker and validation scripts use only the Python
standard library. Measured normalized sizes should be reproduced with
CPython 3.12.14, since \texttt{ast.unparse} is version-dependent.

\begin{thebibliography}{9}
\bibitem{baseline} ProveIt contributors, \emph{Concise presentations of
intuitionistic primitive recursive arithmetic}, repository snapshot,
commit \texttt{9f9545f5ffd9} (full hash on the title page),
12 November 2026.
\href{https://github.com/VladimirReshetnikov/ProveIt/tree/9f9545f5ffd9e0392d4a36639503aca12383450a/Logic/PrimitiveRecursiveArithmetic/ConciseIPRA}{Pinned repository snapshot}.

\bibitem{readme} Same snapshot, \emph{README}, especially ``Further
questions'', item 2, and the raw clause / two-case evaluation entries.
The exact retrieved README is included in the source provenance.

\bibitem{rootfree} Same snapshot,
\texttt{Research/rootfree\_clausal} and
\texttt{Research/reduced\_clause\_calculus}, existing two-case evaluation
and raw clause presentation. The exact matched control path is recorded
in \texttt{PROVENANCE.json}.

\bibitem{evallean} Same snapshot,
\texttt{Lean/ConciseIPRA/NativeEval.lean},
\texttt{NativeClausal.lean} and the cited clause-calculus successors.

\bibitem{biinterp} Same snapshot,
\texttt{Lean/ConciseIPRA/NativeBiInterpretation.lean}.
\href{https://github.com/VladimirReshetnikov/ProveIt/blob/9f9545f5ffd9e0392d4a36639503aca12383450a/Logic/PrimitiveRecursiveArithmetic/ConciseIPRA/Lean/ConciseIPRA/NativeBiInterpretation.lean}{Pinned Lean bi-interpretation endpoint}.

\bibitem{bridge} Same snapshot,
\texttt{Research/certified\_zero/reference/}:
\texttt{inherited\_choice\_basis.tex},
\texttt{inherited\_ordinary\_bridge.tex},
\texttt{scalar\_predecessor.tex},
\texttt{equality\_reflection.tex},
\texttt{primitive\_recursion.tex}.
These are included unchanged under their repository license.
\end{thebibliography}
\end{document}
